\section{Recurrent networks and LSTM (Prof.\ Hochreiter)}

This is Prof.\ Hochreiter's own field: he found the vanishing gradient (1991) and invented LSTM
with Schmidhuber (1997). Expect him to go deep and to enjoy it. Notation of his lecture: input
$\mathbf{x}(t)$, pre-activation $\mathbf{s}(t)$, activation $\mathbf{a}(t)$ or hidden state
$\mathbf{h}(t)$, memory cell $\mathbf{c}(t)$, input weights $W$, recurrent weights $R$, output
weights $V$.

\subsection{Recurrent networks and BPTT}

\begin{summary}{RNN in two minutes}
\textsf{One idea.} An RNN reads a sequence one step at a time and keeps a hidden state: a summary
of everything seen so far. The \emph{same} weights are used at every step, so it works for any
length. Training = backpropagation through time (BPTT): unroll the net over time and backprop
through the unrolled graph.
\[
\mathbf{s}(t)=W\mathbf{x}(t)+R\,\mathbf{a}(t-1),\qquad \mathbf{a}(t)=f(\mathbf{s}(t)),\qquad
\hat{\mathbf{y}}(t)=V\mathbf{a}(t),\qquad L=\sum_{t=1}^{T}L(t).
\]
\[
\bm\delta(t)=\mathrm{diag}\bigl(f'(\mathbf{s}(t))\bigr)\Bigl(V^\top\frac{\partial L(t)}{\partial\hat{\mathbf{y}}(t)}+R^\top\bm\delta(t+1)\Bigr),
\]
\[
\frac{\partial L}{\partial R}=\sum_t\bm\delta(t)\,\mathbf{a}(t-1)^\top,\qquad
\frac{\partial L}{\partial W}=\sum_t\bm\delta(t)\,\mathbf{x}(t)^\top.
\]
\textsf{Key words:} weight sharing in time \textbullet{} Turing complete \textbullet{} BPTT
\textbullet{} truncated BPTT \textbullet{} RTRL \textbullet{} vanishing gradient in time.
\end{summary}

\pic{Reading a novel. You do not remember every word, you keep a running summary in your head
(the hidden state) and update it with each new sentence, using the same reading skill (the same
weights) on every page. BPTT is re-reading the book backwards to find which pages caused your
wrong guess about the murderer.}

\board{$\mathbf{s}(t)=W\mathbf{x}(t)+R\mathbf{a}(t-1)$,\ $\mathbf{a}(t)=f(\mathbf{s}(t))$,\
$\hat{\mathbf{y}}(t)=V\mathbf{a}(t)$;\
$\bm\delta(t)=\mathrm{diag}(f'(\mathbf{s}(t)))\,(V^\top\partial L(t)/\partial\hat{\mathbf{y}}(t)+R^\top\bm\delta(t+1))$;\
$\partial L/\partial R=\sum_t\bm\delta(t)\mathbf{a}(t-1)^\top$.
Note: $\bm\delta(t+1)$, the future, not $t-1$. Gradient uses $\mathbf{a}(t-1)$, not $\mathbf{s}(t-1)$.}

\begin{qabox}[Hochreiter]{Which recurrent architectures do you know?}
\ans Jordan network: feeds the \emph{output} of $t-1$ back as input. Elman network: feeds the
\emph{hidden} state back, originally with fixed recurrent weights and trained locally in time.
Fully recurrent network: trainable recurrent matrix $R$, every hidden unit sees all others.
NARX: shortcuts to outputs further in the past (delays), which fights vanishing gradients only for a
finite number of steps. Time-delay networks (TDNN): a fixed window, basically a 1D convolution.
Then LSTM and GRU with gates.
\follow ``Jordan training trick?'' Teacher forcing: during training feed the true label of $t-1$
instead of the prediction.
\src{LSTM lecture WS24, \texttt{Exam\_LSTM summary.pdf}.}
\end{qabox}

\begin{qabox}[Hochreiter]{BPTT, truncated BPTT and RTRL: compare them.}
\ans BPTT: forward through the whole sequence, then backward; exact gradient; memory grows
linearly with the length $T$; cost $O(T\cdot W)$. Truncated BPTT: backprop only $k$ steps into
the past; cheaper, but the gradient is an approximation and long dependencies are cut. RTRL:
carries the derivative of the state with respect to all weights forward in time, so it is online
and memory does not grow with $T$, but it is very expensive: $O(N^3)$ numbers stored and $O(N^4)$
operations per step for $N$ hidden units.
\formula RTRL recursion:
$\dfrac{\partial\mathbf{s}(t)}{\partial w}=\dfrac{\partial\mathbf{s}(t)}{\partial w}\Big|_{\text{direct}}
+\dfrac{\partial\mathbf{s}(t)}{\partial\mathbf{s}(t-1)}\dfrac{\partial\mathbf{s}(t-1)}{\partial w}$.
\pic{BPTT is watching the whole match recording afterwards. RTRL is a coach who updates his
notes live during the match: no recording needed, but he must write very fast.}
\src{LSTM lecture WS24, BPTT/RTRL exam questions.}
\end{qabox}

\begin{qabox}[Hochreiter]{Why can a plain RNN not learn long-term dependencies?}
\ans The delta is multiplied at every time step by the recurrent Jacobian
$R^\top\mathrm{diag}(f'(\mathbf{s}))$. Over $k$ steps this is a product of $k$ factors, so it
decays (or explodes) exponentially in $k$. After 10--20 steps the error from the past is
practically zero.
\formula
$\Bigl\lVert\dfrac{\partial\mathbf{a}(t)}{\partial\mathbf{a}(t-k)}\Bigr\rVert
=\Bigl\lVert\prod_{\tau=t-k+1}^{t}\mathrm{diag}\bigl(f'(\mathbf{s}(\tau))\bigr)R\Bigr\rVert
\le\bigl(\max|f'|\,\lVert R\rVert\bigr)^{k}.$
\follow ``Gradient clipping?'' Only fixes \emph{exploding}: it rescales the gradient norm to a
threshold, keeps the direction. It does nothing for vanishing.
\end{qabox}

\subsection{LSTM}

\begin{summary}{LSTM in two minutes}
\textsf{One idea.} Give the network a \emph{memory cell} that is updated by \emph{adding}, not by
multiplying through a squashing function. Then the error can flow back through time unchanged:
the \emph{constant error carousel} (CEC). Gates, small sigmoid switches between 0 and 1, decide
what to write (input gate), what to keep (forget gate) and what to show (output gate).
\[
\begin{aligned}
\mathbf{z}(t)&=\tanh\bigl(W_z\mathbf{x}(t)+R_z\mathbf{h}(t-1)+\mathbf{b}_z\bigr)&&\text{cell input}\\
\mathbf{i}(t)&=\sigma\bigl(W_i\mathbf{x}(t)+R_i\mathbf{h}(t-1)+\mathbf{b}_i\bigr)&&\text{input gate}\\
\mathbf{f}(t)&=\sigma\bigl(W_f\mathbf{x}(t)+R_f\mathbf{h}(t-1)+\mathbf{b}_f\bigr)&&\text{forget gate}\\
\mathbf{o}(t)&=\sigma\bigl(W_o\mathbf{x}(t)+R_o\mathbf{h}(t-1)+\mathbf{b}_o\bigr)&&\text{output gate}\\
\mathbf{c}(t)&=\mathbf{f}(t)\odot\mathbf{c}(t-1)+\mathbf{i}(t)\odot\mathbf{z}(t)&&\text{memory update}\\
\mathbf{h}(t)&=\mathbf{o}(t)\odot\tanh\bigl(\mathbf{c}(t)\bigr)&&\text{output}
\end{aligned}
\]
\textsf{Why it works:} $\dfrac{\partial\mathbf{c}(t)}{\partial\mathbf{c}(t-1)}=\mathrm{diag}(\mathbf{f}(t))$,
which is $\approx I$ when the forget gate is open. The original 1997 LSTM had no forget gate:
$\mathbf{c}(t)=\mathbf{c}(t-1)+\mathbf{i}\odot\mathbf{z}$, derivative exactly $I$. The forget gate
was added by Gers, Schmidhuber and Cummins (2000).
\end{summary}

\pic{A conveyor belt (the cell state) running through a factory. Things ride on it without being
changed. Three workers stand next to it: one decides what to put \emph{on} the belt (input gate),
one decides what to throw \emph{off} (forget gate), one decides what to \emph{show} to the next
department (output gate). Because the belt does not squash anything, an item put on at the
beginning is still there at the end: long-term memory.}

\board{$\mathbf{c}(t)=\mathbf{f}\odot\mathbf{c}(t-1)+\mathbf{i}\odot\mathbf{z}$ \quad
$\mathbf{h}(t)=\mathbf{o}\odot\tanh(\mathbf{c}(t))$ \quad
gates $=\sigma(W\mathbf{x}(t)+R\mathbf{h}(t-1)+\mathbf{b})$, cell input $=\tanh(\cdot)$ \quad
$\partial\mathbf{c}(t)/\partial\mathbf{c}(t-1)=\mathrm{diag}(\mathbf{f}(t))\approx I$.\
Draw the cell: belt on top, three $\sigma$ gates, one $\tanh$ in, one $\tanh$ out.}

\begin{qabox}[Hochreiter]{Explain the constant error carousel.}
\ans In the original LSTM the cell connects to itself with weight 1 and a linear activation. So
the local error in the cell is multiplied by exactly 1 at every step: it neither vanishes nor
explodes. The gates protect this carousel: the input gate protects the cell from irrelevant
inputs, the output gate protects other units from irrelevant cell content.
\formula $\mathbf{c}(t)=\mathbf{c}(t-1)+\mathbf{i}(t)\odot\mathbf{z}(t)\ \Rightarrow\
\partial\mathbf{c}(t)/\partial\mathbf{c}(t-1)=I$.
\src{LSTM lecture WS24, CEC slides; Hochreiter \& Schmidhuber (1997).}
\end{qabox}

\begin{qabox}[Hochreiter]{Why was the forget gate added? What about drift?}
\ans Without a forget gate the cell can only accumulate; on long, continuous input streams the
state grows without limit and the cell saturates. The forget gate lets the network reset its
memory. Related: \emph{internal state drift}, when inputs are mostly of one sign the state drifts
away; the remedy in the original paper is to bias the input gates towards zero at the start
(negative bias), so the cell writes only when needed.
\follow ``Ticker steps?'' Letting the LSTM run extra steps after the input has ended, so it can
process what it has stored before it must answer.
\end{qabox}

\begin{qabox}[Hochreiter]{LSTM variants: peepholes, CIFG, GRU.}
\ans Peephole connections let the gates see the cell state $\mathbf{c}(t-1)$ directly.
CIFG (coupled input and forget gate) sets $\mathbf{f}=1-\mathbf{i}$: write only as much as you
forget. GRU (Cho et al.\ 2014) goes further: no separate cell, an update gate that acts as coupled
input and forget gate, and a reset gate that decides how much the old state enters the candidate.
\formula GRU:\ $\mathbf{u}=\sigma(\cdot)$,\ $\mathbf{r}=\sigma(\cdot)$,\
$\tilde{\mathbf{h}}=\tanh\bigl(W\mathbf{x}(t)+R(\mathbf{r}\odot\mathbf{h}(t-1))\bigr)$,\
$\mathbf{h}(t)=(1-\mathbf{u})\odot\mathbf{h}(t-1)+\mathbf{u}\odot\tilde{\mathbf{h}}$.
\follow Other variants from the lecture: multidimensional LSTM (e.g.\ 2D images, several time
axes), ConvLSTM (convolutions inside the gates, for video).
\end{qabox}

\begin{qabox}[Hochreiter]{Dropout in LSTMs: why is it special?}
\ans Naive dropout on the recurrent state at every step multiplies survival probabilities over
time, so a memory would survive with probability $(1-p)^T$: long-term memory is destroyed. So we
use special forms: dropout only on the non-recurrent connections, variational dropout (the same
mask at every time step), or zoneout (randomly \emph{keep} the previous state instead of zeroing
units). Dropout is switched off at evaluation.
\end{qabox}

\begin{qabox}[Hochreiter, Schl\"uter]{Bidirectional LSTM, and one label for a whole sequence.}
\ans A bidirectional LSTM runs one LSTM forward and one backward and concatenates the two hidden
states, so each step sees past and future; it needs the full sequence, so not for real-time use.
For one label per sequence use the last hidden state or global average pooling over all states,
then a classifier. Seq2seq: an encoder LSTM reads the input, a decoder LSTM writes the output.
\end{qabox}

\subsection{Beyond LSTM: xLSTM, Mamba, Hopfield}

\begin{summary}{Modern sequence models from Linz and elsewhere}
\textsf{xLSTM} (Beck, Hochreiter et al.\ 2024): (1) \emph{exponential gating} with a normalizer
state, so the model can revise a storage decision (a new, more important input can overwrite the
old one); (2) \textsf{sLSTM}: scalar memory with memory mixing and a stabilizer against overflow;
(3) \textsf{mLSTM}: a \emph{matrix} memory $C\in\mathbb{R}^{d\times d}$ updated by a covariance rule,
fully parallelizable; its size does not grow with the sequence length.
\[
\text{mLSTM: }C_t=f_tC_{t-1}+i_t\,\mathbf{v}_t\mathbf{k}_t^\top,\quad
\mathbf{n}_t=f_t\mathbf{n}_{t-1}+i_t\mathbf{k}_t,\quad
\mathbf{h}_t=\mathbf{o}_t\odot\frac{C_t\mathbf{q}_t}{\max(|\mathbf{n}_t^\top\mathbf{q}_t|,1)},\quad i_t=\exp(\tilde i_t).
\]
\textsf{Mamba}: a \emph{selective} state space model, $\mathbf{h}_t=\bar A\mathbf{h}_{t-1}+\bar B\mathbf{x}_t$,
$\mathbf{y}_t=C\mathbf{h}_t$, where $B$, $C$ and the step $\Delta$ depend on the input; linear in
sequence length. \
\textsf{Modern Hopfield networks} (Ramsauer, Hochreiter et al.\ 2020): the update
$\bm\xi^{\text{new}}=X\,\mathrm{softmax}(\beta X^\top\bm\xi)$ is exactly Transformer attention.
\end{summary}

\pic{Old LSTM memory is a notebook where you write in pencil and can only fade old notes. xLSTM's
exponential gate is a highlighter: a new, important note can outshine everything before it. The
mLSTM matrix memory is a filing cabinet: store a value under a key ($\mathbf{v}\mathbf{k}^\top$),
retrieve it later with a query ($C\mathbf{q}$).}

\begin{qabox}[Hochreiter]{What is new in xLSTM, and why?}
\ans LSTM has three limits: it cannot revise a storage decision, its memory is a small scalar
cell, and it cannot be parallelized over time because of memory mixing. xLSTM answers each:
exponential gating with normalization revises storage decisions; mLSTM gives a matrix memory
with key--value storage; and mLSTM drops memory mixing, so it trains in parallel like a
Transformer while inference stays linear in the sequence length.
\formula $C_t=f_tC_{t-1}+i_t\mathbf{v}_t\mathbf{k}_t^\top$, retrieval $C_t\mathbf{q}_t$; normalizer
$\mathbf{n}_t=f_t\mathbf{n}_{t-1}+i_t\mathbf{k}_t$.
\src{LSTM lecture WS24, xLSTM exam questions.}
\end{qabox}

\begin{qabox}[Hochreiter]{What connects Hopfield networks and attention?}
\ans A modern Hopfield network stores patterns $X$ and retrieves the one closest to a query
$\bm\xi$ with $\bm\xi^{\text{new}}=X\,\mathrm{softmax}(\beta X^\top\bm\xi)$. With $\beta=1/\sqrt{d_k}$
and learned projections this is exactly the attention formula. So attention is associative
memory retrieval, and the network can store exponentially many patterns.
\end{qabox}
